
\begin{proposition}[Proved equality comparison]\label{cdp:st:C.2.VII.4}

Retain the irreducible class VII situation of
\ref{cdp:st:C.2.VII.3}. Let $C=\Gamma_{\rm red}$, let $n_0$
be the number of its rational-normalized components,
and put $v(C)=\sum_x(s_x-1)$, where $s_x$ is the
number of branches at a singular point of $C$.
On the minimal resolution $Z$ let $D$ be the
resolution conductor, let $r$ count the rational
curves in $D_{\rm red}$, and let $a$ count the
zero-coefficient exceptional curves lying over
points outside the conductor. Let $b$ count the
positive-coefficient exceptional fibres over
such points, as in \ref{cdp:st:C.2.VII.3}. Then $e(S)=0$
is equivalent to the following explicit
normalization and curve data:
\[
 b=0,\qquad b_2(Z)=r+a,
 \qquad
 \begin{cases}
 n_0=0, v(C)=0,\quad\text{or}\\
 n_0=1, v(C)=2.
 \end{cases}
 \tag{C.VII.9}\label{cdp:eq:C.VII.9}
\]
In the first curve alternative all components
have elliptic normalization and the sum of
their delta invariants is at most one. In the
second the total delta invariant is exactly
two, and every singular point satisfies
$\delta_x=s_x-1$.

If the actual homomorphism
$\Pic^0(Y)\to\Pic^0(Z)$ is nonconstant,
there is the sharper geometric equivalence:
$e(S)=0$ exactly when its normalization is a
smooth minimal Hopf surface and the reduced
conductor image and inverse image are disjoint
unions of smooth elliptic curves, with finite
unramified maps between them.
\end{proposition}
\begin{proof}
The exact comparison already proved in
\ref{cdp:st:C.2.VII.3} is
\[
 e(S)=b_2(Z)-r-a+b+e(C),
 \qquad b_2(Z)-r-a\ge0,\quad b\ge0,\quad e(C)\ge0.
 \tag{C.VII.10}\label{cdp:eq:C.VII.10}
\]
Every normalized component of $C$ has genus
zero or one. Its topological normalization
formula is $e(C)=2n_0-v(C)$. The stronger
bound in \ref{cdp:st:C.2.VII.2} is $e(C)\ge n_0-1$.
Thus $e(C)=0$ permits precisely $n_0=0,1$,
with $v(C)=0,2$, respectively. This proves
both directions of \eqref{cdp:eq:C.VII.9}, using actual
branch numbers rather than an assumption
about a local singularity type.
The same proof gives
$\chi(\OO_C)=n_0-\sum\delta_x\ge-1$.
For $n_0=0$ this bounds the total delta
invariant by one. For $n_0=1$, the value
$v(C)=2$ and $\delta_x\ge s_x-1$ force
$\sum\delta_x=2$ and equality at every point.
These are exact consequences for fibres which
satisfy the criterion; no separate claim
that all such abstract curve patterns admit
a torus smoothing is used.

Now suppose that the Picard pullback is
nonconstant. The argument of \ref{cdp:st:A.5.12} applies
to $\Pic^0(Y)$ even with only rational
$H^2(Y,\Q)=0$. Namely its nonzero differential
is the actual composite
\[
 H^1(Y,\OO_Y)\longrightarrow H^1(S,\OO_S)
 \longrightarrow H^1(T,\OO_T)
 \longrightarrow H^1(Z,\OO_Z).
 \tag{C.VII.11}\label{cdp:eq:C.VII.11}
\]
The middle arrow consequently has positive
rank. Put $\varepsilon=h^0(S,\omega_S)$;
the torsion canonical relation makes
$\varepsilon\in\{0,1\}$ and
$h^1(S,\OO_S)=1+\varepsilon$.
The normalization sequence then gives
$h^0(\mathcal Q)\le\varepsilon$.
Conductor duality and restriction onto $C$
therefore give
\[
 h^1(C,\omega_S|_C)\le\varepsilon.
 \tag{C.VII.12}\label{cdp:eq:C.VII.12}
\]
If $\varepsilon=0$, degree-zero
Riemann--Roch gives $\chi(\OO_C)\ge0$.
If $\varepsilon=1$, the torsion line
$\omega_S$ is trivial, and
$h^0(C,\OO_C)\ge1$ gives the same inequality.
The exact formula (C.INT.26) now implies
that $e(C)=0$ holds precisely when $C$ is
a disjoint union of smooth elliptic curves.

Assume $e(S)=0$. Every rational curve in
$D_{\rm red}$ must be resolution-exceptional,
because its image lies in $C$ and cannot
dominate an elliptic curve. A minimal class
VII model with positive second Betti number
and a torsion anticanonical divisor is one
of the proved configurations used in
\ref{cdp:st:C.2.VII.1}. An intermediate or
Inoue--Hirzebruch model has no elliptic
curves, and therefore could not dominate
the nonempty elliptic conductor. A
parabolic Inoue model has a rational cycle
of square zero and exactly $b_2$ rational
curves. Its point blowups therefore still
have exactly $b_2(Z)$ rational curves.
The equality $b_2(Z)=r+a$ makes every one
of them resolution-exceptional, including
any zero-coefficient blowup curves. Thus
the full total transform of the minimal
rational cycle, whose square remains zero,
would be exceptional. This contradicts
negative definiteness of a resolution's
exceptional intersection matrix.
Consequently the minimal model is Hopf.
All rational curves on $Z$ are now the
exceptional curves of its point blowups,
and their number is $b_2(Z)$.
The equality $b_2(Z)=r+a$ makes all of
them resolution-exceptional. As in the
last paragraph of \ref{cdp:st:C.2.INT.3}, a nonempty
blowup sequence would have a last
$(-1)$ curve, contradicting minimality
of the resolution. Thus the normalization
is already a smooth minimal Hopf surface.
Its reduced conductor curves are elliptic;
their finite maps to the smooth elliptic
curves of $C$ are unramified by
Riemann--Hurwitz.
Conversely those normalization and
conductor data have Euler number zero
by the exact normalization formula.
This proves the sharper equivalence.
\end{proof}


